\chapter{Introduction}
\label{cha:intro}

Even though most of the common graph theory notions and objects are well
understood, the area lacks a consistent and unified terminology. For that
matter, we devote this chapter to the introduction of the basic definitions
along with notation. We also present some ground results on which we build our
work. Even for the familiarized reader, we encourage to at least skim through
the sections, as we introduce some non-standard notations for operations between
graphs, as well as conventions for certain types of partitions of graphs. All
definitions not considered in here can be found in either \cite{bondy2008} or
\cite{diestel2017}.

%%%%%%%%%%%%%%%%%%%%%%%%%%%%%%%%%%%%%%%%%%%%%%%%%%%%%%%%%%%%%%%%%%%%%%%%%%%%%%%%
%%%%%%%%%%%%%%%%%%%%%%%%%%%%%%%%%%%%%%%%%%%%%%%%%%%%%%%%%%%%%%%%%%%%%%%%%%%%%%%%
%%%%%%%%%%%%%%%%%%%%%%%%%%%%%%%%%%%%%%%%%%%%%%%%%%%%%%%%%%%%%%%%%%%%%%%%%%%%%%%%
%%%%%%%%%%%%%%%%%%%%%%%%%%%%%%%%%%%%%%%%%%%%%%%%%%%%%%%%%%%%%%%%%%%%%%%%%%%%%%%%
\section{Graphs and algorithms}
\label{sec:graphs}

All graphs considered in this work are finite, undirected with neither loops nor
parallel edges. We denote by $V$ and $E$ the vertex and edge sets of a graph
$G$, respectively. For a vertex $v\in V$, $N(v)$ is the set of neighbors of $v$,
and the degree of a vertex is $d(v) = |N(v)|$. For $S\subseteq V$, we write
$G[S]$ for the subgraph of $G$ induced by $S$. When two graphs $G$ and $H$ are
isomorphic, we write $G \cong H$.

The complement of a graph $G$ is denoted by $\overline{G}$, and the disjoint
union of two graphs $G$ and $H$ is referred to as $G + H$. For a positive
integer $k$, let $kG$ denote the disjoint union of $k$ copies of $G$. We denote
by $G-S$ the subgraph $G[V\setminus S]$ --- if $S$ is a singleton, then we write
$G-v$ instead of $G-\l{v}$.

For a positive integer $n$, we denote by $K_n$, $C_n$ and $P_n$ the complete
graph, the cycle, and the path on $n$ vertices, respectively. We say that a set
$S\subseteq V$ is independent or stable if $G[S] \cong \overline{K_{|S|}}$, and
$S$ is a clique if $G[S] \cong K_{|S|}$. We say that a graph is a cluster if it
is the disjoint union of complete graphs, and a graph is a complete multipartite
if it is the complement of a cluster. We call a graph $G$ chordal if it does not
contain any induced cycle of length at least four, that is, if $C$ is a cycle in
$G$ of length at least four, then there is a pair of adjacent non-consecutive
vertices of $C$; such an edge is called a chord of $C$.

We say that a vertex $v$ is universal if it is adjacent to all other vertices of
$G$, and isolated if it is adjacent to none. Additionally, $v$ is said to be
simplicial if $N(v)$ is a clique.

Finally, in \Cref{fig:classic} we present two classic graphs that we will
frequently refer to.

\begin{figure}[htb!]
    \centering
    \begin{tikzpicture}
        \matrix[column sep=1.5cm]{
            \node {
                \begin{tikzpicture}[every node/.style={vertex}, edge]
                    \foreach \i in {0,...,3}
                        \node (\i) at ({(360/4)*\i + 90}:1) {};

                    \draw (0) to (1);
                    \draw (2) to (1);
                    \draw (2) to (3);
                    \draw (0) to (3);
                    \draw (0) to (2);
                \end{tikzpicture}
            }; &
            \node {
                \begin{tikzpicture}[every node/.style={vertex}, edge]
                    \foreach \i in {0,...,3}
                        \node (\i) at ({(360/4)*\i + 45}:1) {};
                    \path let \p1=(0) in node (r) at (0,2*\y1) {};

                    \draw (0) to (1);
                    \draw (1) to (2);
                    \draw (2) to (3);
                    \draw (3) to (0);
                    \draw (r) to (0);
                    \draw (1) to (r);
                \end{tikzpicture}
            }; \\
            \node {diamond}; &
            \node {house}; \\
        };
    \end{tikzpicture}
    \caption{Some classic graphs.}
    \label{fig:classic}
\end{figure}

Now we briefly discuss some notions regarding algorithms on graphs. Mainly, we
are interested in the running time of algorithms. For that matter, we use the
standard big-$O$ notation to describe asymptotic behavior. For all algorithms
described in this work, whenever we say that it runs in linear time, we mean
that its time complexity is $O(|V| + |E|)$.

It is well known that for an input graph $G$, it is possible to order its
adjacency lists to have compatible orderings, i.e., there is a linear ordering
of $V$ such that within each adjacency list, the vertices appear in that
ordering.   Therefore, we will assume that this is always the case.

Given a partition $\mathcal{P}$ of constant size for an input graph $G$, it is
possible to construct copies of the subgraphs $G[P]$ for each $P \in
\mathcal{P}$ in linear time by traversing the adjacency lists of each vertex in
$G$.   Thus, given an input graph $G$, it is possible to execute a linear-time
algorithm in each $G[P]$ in time proportional to the number of vertices and
edges of $G[P]$.   It is of particular interest for this work to know that
$G[P]$ can be tested to be split, empty, complete or a cluster graph in linear
time.

, such a thing can be accomplished by
running a BFS to find all connected components, and then verifying if the degree
of each vertex is equal to the size of its component minus one; otherwise, a
second pass of BFS can be executed to find an induced copy of $P_3$.


%%%%%%%%%%%%%%%%%%%%%%%%%%%%%%%%%%%%%%%%%%%%%%%%%%%%%%%%%%%%%%%%%%%%%%%%%%%%%%%%
%%%%%%%%%%%%%%%%%%%%%%%%%%%%%%%%%%%%%%%%%%%%%%%%%%%%%%%%%%%%%%%%%%%%%%%%%%%%%%%%
%%%%%%%%%%%%%%%%%%%%%%%%%%%%%%%%%%%%%%%%%%%%%%%%%%%%%%%%%%%%%%%%%%%%%%%%%%%%%%%%
%%%%%%%%%%%%%%%%%%%%%%%%%%%%%%%%%%%%%%%%%%%%%%%%%%%%%%%%%%%%%%%%%%%%%%%%%%%%%%%%
\section{Exact neighborhoods}
\label{sec:exact-neigh}

%TODO: Definiciones de vecindades exactas y breve explicacion de como calcularlas en tiempo lineal.



%%%%%%%%%%%%%%%%%%%%%%%%%%%%%%%%%%%%%%%%%%%%%%%%%%%%%%%%%%%%%%%%%%%%%%%%%%%%%%%%
%%%%%%%%%%%%%%%%%%%%%%%%%%%%%%%%%%%%%%%%%%%%%%%%%%%%%%%%%%%%%%%%%%%%%%%%%%%%%%%%
%%%%%%%%%%%%%%%%%%%%%%%%%%%%%%%%%%%%%%%%%%%%%%%%%%%%%%%%%%%%%%%%%%%%%%%%%%%%%%%%
%%%%%%%%%%%%%%%%%%%%%%%%%%%%%%%%%%%%%%%%%%%%%%%%%%%%%%%%%%%%%%%%%%%%%%%%%%%%%%%%
\section{Hereditary properties}
\label{sec:hereditary}

%TODO: Obstrucciones, libres de F, consecuencias basicas de la definicion.

%%%%%%%%%%%%%%%%%%%%%%%%%%%%%%%%%%%%%%%%%%%%%%%%%%%%%%%%%%%%%%%%%%%%%%%%%%%%%%%%
%%%%%%%%%%%%%%%%%%%%%%%%%%%%%%%%%%%%%%%%%%%%%%%%%%%%%%%%%%%%%%%%%%%%%%%%%%%%%%%%
%%%%%%%%%%%%%%%%%%%%%%%%%%%%%%%%%%%%%%%%%%%%%%%%%%%%%%%%%%%%%%%%%%%%%%%%%%%%%%%%
%%%%%%%%%%%%%%%%%%%%%%%%%%%%%%%%%%%%%%%%%%%%%%%%%%%%%%%%%%%%%%%%%%%%%%%%%%%%%%%%
\section{Chordal graphs and its subclasses}
\label{sec:chordal-defs}

A well-studied class of graphs is that of chordal graphs. In
\cite{golumbic2004}, Golumbic presents some characterizations of chordal graphs,
one of which yields a linear-time algorithm to recognize them. 

%TODO: Definicion cordal, caracterizaciones relevantes (peo), linealidad de clan maximo y similares.

%TODO: Definicion split y caracterizaciones relevantes. Reconocimiento en tiempo
%lineal (sucesion de grados y algoritmo certificador). Partición con K vacío o
%clan maximal no trivial.


%%%%%%%%%%%%%%%%%%%%%%%%%%%%%%%%%%%%%%%%%%%%%%%%%%%%%%%%%%%%%%%%%%%%%%%%%%%%%%%%
%%%%%%%%%%%%%%%%%%%%%%%%%%%%%%%%%%%%%%%%%%%%%%%%%%%%%%%%%%%%%%%%%%%%%%%%%%%%%%%%
%%%%%%%%%%%%%%%%%%%%%%%%%%%%%%%%%%%%%%%%%%%%%%%%%%%%%%%%%%%%%%%%%%%%%%%%%%%%%%%%
%%%%%%%%%%%%%%%%%%%%%%%%%%%%%%%%%%%%%%%%%%%%%%%%%%%%%%%%%%%%%%%%%%%%%%%%%%%%%%%%
\section{Polarity}
\label{sec:polarity-defs}

\section{Monopolarity}
\label{cap:mono}

\section{Previous results}
\label{sec:previous}

\begin{proposition}
    Let $G$ and $H$ be graphs.   If $G$ and $H$ are monopolar, then $G+H$ is
    monopolar.
\end{proposition}
\begin{proof}
    Let $(A_1, B_1)$ and $(A_2, B_2)$ be the monopolar partitions of $G$ and
    $H$, respectively.   Since $A_1$ and $A_2$ are independent, we have that
    $A_1 \cup A_2$ is also independent.   Analogously, $G[B_1]$ and $G[B_2]$ are
    clusters, so their union is also a cluster.   Hence, $(A_1 \cup A_2, B_1
    \cup B_2)$ is a monopolar partition of $G+H$.
    
\end{proof}

\begin{corollary}
    A graph is monopolar if and only if each of its connected components is
    monopolar.
\end{corollary}

\begin{corollary}
    Every minimal obstruction for monopolarity is a connected graph.
\end{corollary}


%TODO: Polarity definitions and basic results (NP-completness, etc.)